\section{Absolute value, squares, and intervals}\label{sec:magnitude}
The \emph{absolute value} $|x|$ of a real number $x$ is its distance from zero on the number line. Thus $|3|=|-3|=3$ and $|0|=0$. As a formula,
\[
|x|=\begin{cases}x&\text{if }x\ge0,\\-x&\text{if }x<0.\end{cases}
\]
In the second case the result is positive, since the negative of a negative number is positive: $|-3|=-(-3)=3$. So $|x|\ge0$ for every $x$. Section~\ref{ch01:sec:notation} found that $\sqrt{x^2}$ is $x$ when $x\ge0$ and $-x$ when $x<0$, so $\sqrt{x^2}=|x|$.

For all real numbers $c$ and $x$,
\[
|cx|=|c|\,|x|.
\]
We check this case by case, according to the signs of $c$ and $x$.
\begin{itemize}
\item If $c=0$ or $x=0$, then $cx=0$, and both sides are zero.
\item If $c>0$ and $x>0$, then $cx>0$, and both sides equal $cx$.
\item If $c<0$ and $x<0$, then $cx>0$, so $|cx|=cx$, while $|c|\,|x|=(-c)(-x)=cx$.
\item If one of $c$ and $x$ is positive and the other negative, then $cx<0$, so $|cx|=-cx$, while $|c|\,|x|$ is $c(-x)$ if $c>0$ and $(-c)x$ if $c<0$, and both equal $-cx$.
\end{itemize}
Taking $c=-1$ gives $|-x|=|x|$: changing the sign of a number does not change its distance from zero.

For a number $b\ge0$, the statement $|x|\le b$ says that $x$ is at distance at most $b$ from zero:
\[
|x|\le b\quad\Longleftrightarrow\quad -b\le x\le b.
\]
To check this, use the two cases in the formula for $|x|$. If $x\ge0$, then $|x|\le b$ says $x\le b$, and $-b\le x$ holds automatically, since $-b\le0\le x$. If $x<0$, then $|x|\le b$ says $-x\le b$, which becomes $x\ge-b$ after multiplying by $-1$, and $x\le b$ holds automatically, since $x<0\le b$. The same argument with strict inequalities shows that, for $b>0$,
\[
|x|<b\quad\Longleftrightarrow\quad -b<x<b.
\]
Both rules apply to a difference $x-c$ in place of $x$. The number $|x-c|$ is the distance between $x$ and $c$ on the number line; for instance, $|7-5|=|5-7|=2$. For $b>0$, the second rule turns $|x-c|<b$ into $-b<x-c<b$, and adding $c$ to each part gives $c-b<x<c+b$: the number $x$ lies less than $b$ away from $c$. For example, $|x-5|<1$ means $4<x<6$.

Now take $b=2$: $|x|\le2$ is the conjunction ``$-2\le x$ and $x\le2$.'' By De Morgan's law, its negation $|x|>2$ is the disjunction of the negations of the two parts, ``$x<-2$ or $x>2$'': a condition with two separate pieces, one on each side of zero. The same two-sided form appears when an inequality involves a square.

We need two facts about squares. First, squaring preserves order for nonnegative numbers. Suppose $0\le a<b$. Multiplying out gives
\[
(b-a)(b+a)=b^2+ba-ab-a^2=b^2-a^2.
\]
Both factors on the left are positive: $b-a>0$ because $a<b$, and $b+a>0$ because $b>a\ge0$. So $b^2-a^2>0$, that is, $a^2<b^2$. In the same way, if $0\le a\le b$, then both factors are nonnegative, so $a^2\le b^2$. The condition that the numbers are nonnegative cannot be dropped: $-3<2$, but $(-3)^2=9$ is larger than $2^2=4$.

Second, $x^2=|x|^2$ for every real number $x$. If $x\ge0$, this says $x^2=x^2$. If $x<0$, it says $x^2=(-x)^2$, which is true because $(-x)(-x)=x\cdot x$.

Now we can solve the inequality $x^2>4$ completely, by comparing $|x|$ with $2$.
\begin{itemize}
\item If $|x|>2$, the squaring rule for $<$ with $a=2$ and $b=|x|$ gives $4<|x|^2=x^2$.
\item If $|x|\le2$, the squaring rule for $\le$ with $a=|x|$ and $b=2$, allowed because $|x|\ge0$, gives $x^2=|x|^2\le4$.
\end{itemize}
The first case shows that $|x|>2$ implies $x^2>4$. The second shows that $|x|\le2$ implies $x^2\le4$, which is the contrapositive of ``$x^2>4$ implies $|x|>2$.'' Together they prove
\[
x^2>4\quad\Longleftrightarrow\quad|x|>2\quad\Longleftrightarrow\quad x<-2\ \text{or}\ x>2.
\]
For instance, $x=-3$ satisfies $x^2=9>4$ even though $x$ is not greater than $2$. The tempting shortcut ``take square roots to get $x>2$'' loses the whole negative half of the answer.

Ranges of real numbers like these have a standard notation. For real numbers $a<b$, each of the following sets is called an \emph{interval}:
\begin{itemize}
\item $[a,b]$ is the set of all real $x$ with $a\le x\le b$;
\item $(a,b)$ is the set of all real $x$ with $a<x<b$;
\item $[a,b)$ is the set of all real $x$ with $a\le x<b$, and $(a,b]$ is the set of all real $x$ with $a<x\le b$.
\end{itemize}
A square bracket includes its endpoint, and a round bracket excludes it: $1$ and $3$ belong to $[1,3]$ but not to $(1,3)$. An interval may also have no endpoint on one side: $[0,\infty)$ is the set of real $x$ with $x\ge0$, and $(-\infty,-2)$ the set of real $x$ with $x<-2$. The symbol $\infty$, read ``infinity,'' is not a real number; it only says that there is no endpoint on that side, so it never goes next to a square bracket. In this notation, $|x|\le2$ says that $x$ belongs to $[-2,2]$, and $x^2>4$ says that $x$ belongs to $(-\infty,-2)$ or to $(2,\infty)$. Round brackets are also used for an \emph{ordered pair}, such as $(1,3)$, which lists two entries in a definite order (Chapter~\ref{ch:maps}); the surrounding sentence tells you which is meant.

\begin{principle}{Habit 2: keep a statement contract}
Before working on a statement, write down its \emph{contract}: the domain of each variable, the hypotheses, the conclusion, and the order in which the quantities are chosen. Whenever the statement changes, for example when you or an assistant repair or generalize it, compare the new version with the contract: did a hypothesis or the conclusion become weaker or stronger? A changed statement can be true and still fail to answer the original question. In Section~\ref{ch01:sec:pattern}, an assistant that replaced ``integers'' by ``positive integers'' gave a correct proof of a weaker statement, one that says nothing about a sum such as $-3+5$.
\end{principle}
